\scanpage{68}
\mspage{12}

\deleted{\deleted{\unsure{tel qu'il y|tel qu'il est}} compatible avec les localisateurs, et tel qu'il existe un foncteur $g:C\unsure{\text{---}|\longrightarrow}A$ compatible avec les localisateurs \deleted{\unsure{et tel que}} avec $fg=p$, $f$ induit une équivalence \deleted{\unsure{sur}} \unsure{entre} les catégories localisées}
\[ \deleted{A\Sigma^{-1}\xrightarrow{\approx} C'{W'}^{-1}.} \]
\diagram{\deleted{\begin{tikzcd}[ampersand replacement=\&,column sep=small]
 C\arrow[rr,"g"]\arrow[dr,"p"'] \& \& A\arrow[dl,"f"]\\
 \& C'\arrow[ul,dashed,bend left=42,"q"']\arrow[ur,dashed,bend right=42,"f'"'] \&
\end{tikzcd}}}

\deleted{\uline{Dém} On va définir un foncteur en sens inverse de $f$}
\[ \deleted{f':C'\longrightarrow A,} \]
\deleted{savoir $f'=gq$. Je dis que le couple $(f,f')$ satisfait aux \unsure{m\^{e}mes} hypothèses que le couple $(p,q)$. En effet}
\[ \deleted{f'f=gqf} \]

Alors $p$ et $\unsure{\varphi|q}$ induisent des \deleted{\unsure{quasi}} équivalences quasi-inverses l'une de l'autre en passant aux catégories localisées (elles n'ont pas de \unsure{mérite}).

De plus, pour \unsure{tte|tt} sous-catégorie pleine $C'_0$ de $C'$, telle que \deleted{$\unsure{C}$} $p$ se factorise \inserted{(en $p_0:C\unsure{\text{---}|\to}C'_0$)} par $C'_0$, si on désigne par $W'_0$ le localisateur induit \deleted{\unsure{dans}} $C'_0$ par $C'$, \deleted{\unsure{les de}} alors $p_0:C\to C'_0$ et $\varphi_0=\varphi\vert C'_0:C'\unsure{\text{---}|\to}C$ induisent également des équivalences qu. inv. l'une de l'autre sur les localisés.

En effet, $\varphi_0p_0=(\varphi i)p_0=\varphi(ip_0)=\varphi p\overset{u}{\text{---}}\mathrm{id}_C$, (ici $i:C'_0\hookrightarrow C'$ désigne l'inclusion), et $\deleted{\unsure{g}}F\overset{\mathrm{d\acute ef}}{=}p_0\varphi_0=p_0\varphi i$ satisfait à
\[ iF=\underbrace{ip_0}_{p}\varphi i=(p\varphi)i. \]

Pour définir \deleted{$\unsure{i}$} $F\overset{v_0}{\text{---}}\mathrm{id}_{C'_0}$, qui est dans $W'_0$ arg. par argument

% archive sheet number: 67
